\documentclass[12pt]{article}
\usepackage[T1]{fontenc}
\usepackage[utf8]{inputenc}
\usepackage{lmodern}
\usepackage{textcomp}
\usepackage{enumitem}
\usepackage{geometry}
\usepackage{graphicx}
\usepackage{amsmath, amssymb}
\geometry{margin=1in}
\begin{document}
\begin{center}
Mathematics - Undergraduate Level Assessment 11 \
Total Marks: 40 \
Answer all questions.
\end{center}

\section*{Multiple Choice (10 marks)}
\begin{enumerate}[label=\arabic*.]
\item Statement 1 | If H is a subgroup of a group G and a belongs to G, then aH = Ha. Statement 2 | If H is normal of G and a belongs to G, then ah = ha for all h in H.
\begin{enumerate}[label=(\alph*)]
    \item True, True
    \item False, False
    \item True, False
    \item False, True
\end{enumerate}
\item Compute the product in the given ring. (12)(16) in Z\_24
\begin{enumerate}[label=(\alph*)]
    \item 0
    \item 1
    \item 4
    \item 6
\end{enumerate}
\item For T: Z x Z -\textgreater{} Z where T(1, 0) = 3 and T(0, 1) = -5, find T(-3,2).
\begin{enumerate}[label=(\alph*)]
    \item -19
    \item -10
    \item 19
    \item 10
\end{enumerate}
\item Find the characteristic of the ring Z\_3 x Z\_3.
\begin{enumerate}[label=(\alph*)]
    \item 0
    \item 3
    \item 12
    \item 30
\end{enumerate}
\item Statement 1 | Any set of two vectors in $R^2$ is linearly independent. Statement 2 | If V = span(v 1, ... , vk) and \{v 1, ... , vk\} are linearly independent, then dim(V) = k.
\begin{enumerate}[label=(\alph*)]
    \item True, True
    \item False, False
    \item True, False
    \item False, True
\end{enumerate}
\end{enumerate}

\section*{True / False (10 marks)}
\textit{Answer True or False}
\begin{enumerate}[label=\arabic*.]
\setcounter{enumi}{5}
\item True or False: The element 4 is a generator for the finite field Z\_7.
\item True or False: If 3 x + 7 y is divisible by 11, then 9 x + 4 y must also be divisible by 11 for positive integers x and y.
\item True or False: In the group G = \{2, 4, 6, 8\} under multiplication modulo 10, the identity element is 6.
\item True or False: It is possible to have a nontrivial homomorphism from a finite group into an infinite group.
\item True or False: If g is a group element and $g^n$ = e, then the order of g is always equal to n regardless of lower powers.
\end{enumerate}

\section*{Long Form Response (20 marks)}
\textit{Answer each question with a clear and complete explanation.}
\begin{enumerate}[label=\arabic*.]
\setcounter{enumi}{10}
\item A circle centered at $A$ with a radius of $1$ and a circle centered at $B$ with a radius of $4$ are externally tangent. A third circle is tangent to the first two and to one of their common external tangents as shown. What is the radius of the third circle? [asy] draw((-3,0)--(7.5,0)); draw(Circle((-1,1),1),linewidth(0.7)); draw(Circle((3,4),4),linewidth(0.7)); draw(Circle((0.33,0.44),0.44),linewidth(0.7)); dot((-1,1)); dot((3,4)); draw((-1,1)--(-2,1)); draw((3,4)--(7,4)); label("$A$",(-1,1),E); label("$B$",(3,4),W); label("1",(-1.5,1),N); label("4",(5,4),N); [/asy]
\item Discuss how to determine the maximum area of a triangle formed with one vertex at the center of a circle of radius 1 and the other two vertices on the circle.
\end{enumerate}

\vfill
\noindent\textit{End of Paper}
\end{document}